\documentclass[11pt]{article}
\usepackage{amsmath,amssymb,amsthm}
\usepackage{booktabs}
\usepackage{enumitem}
\usepackage{hyperref}
\usepackage[margin=1in]{geometry}

\hypersetup{colorlinks=true,linkcolor=black,citecolor=blue,urlcolor=blue}

\newtheorem{theorem}{Theorem}
\newtheorem{lemma}[theorem]{Lemma}
\newtheorem{corollary}[theorem]{Corollary}
\newtheorem{definition}{Definition}
\newtheorem{remark}{Remark}
\newtheorem{proposition}[theorem]{Proposition}

\newcommand{\amt}{\mathrm{amt}}
\newcommand{\spread}{\sigma}
\newcommand{\bal}{\mathrm{bal}}

\title{Lux Adopts the Liquidity Protocol\\
{\large Composition of Quasar 3.0 PQ Finality with the Oracle-Mirrored
AMM Settlement Layer}}
\author{Zach Kelling\\
\textit{Lux Industries, Inc.}\\
\texttt{research@lux.network}}
\date{Adoption: 2026-04-20}

\begin{document}
\maketitle

\begin{abstract}
On 2026-04-20, the Lux Network formally adopts the
\emph{Liquidity Protocol} (Liquidity.io, launched 2026-04-01) as a
canonical regulated digital-securities settlement layer over the Lux
C-Chain (EVM). This note proves three results: (i)~every
Liquidity-Protocol-compliant transaction settled on a Lux execution
chain inherits Quasar~3.0 post-quantum finality;
(ii)~the Liquidity Protocol's per-trade atomicity invariant composes
with QuasarSTM 4.0 serialisability (LP-135) without weakening either;
(iii)~the cross-chain replay protection of the Liquidity Protocol's
oracle-mirrored AMM (OMA) is implied by QuasarCert subject-binding
(Theorem~7.8 of \cite{quasar-cert}).  We rely only on results that are
already machine-checked in \texttt{luxfi/proofs} and on the published
Liquidity Protocol specification.  Adoption is forward-only:
2026-04-21 onwards is reserved for QuasarSTM~4.0 production
activation~(LP-135) and for D-Chain landing per LP-134.
\end{abstract}

\tableofcontents

%% ==================================================================
\section{Background and adoption boundary}
\label{sec:background}
%% ==================================================================

The Lux Network operates a multi-chain topology specified in
LP-134~\cite{lp134}: P-, C-, X-, Q-, Z-, A-, B-, M-, F-Chain
(the activation set) plus a forthcoming D-Chain reserved for the
Liquidity stack. Consensus across all chains is Quasar~3.0
(LP-020, LP-105). The Liquidity Protocol~\cite{lp-paper} is a
self-contained settlement protocol designed for any EVM-compatible
chain. It pairs an Oracle-Mirrored AMM (OMA) contract with an
Alpaca-backed Liquidity ATS brokerage account, settles atomically, and
reverts on insufficient inventory.

Until 2026-04-20 the Lux Network referenced the Liquidity Protocol as a
counterparty integration only.  This document upgrades that status to
\emph{adoption}: the protocol is recognised as a canonical
digital-securities lane on Lux and its proofs become an extension of
the Lux proof corpus.

\begin{definition}[Adoption boundary]
\label{def:boundary}
``Lux adopts the Liquidity Protocol'' means:
\begin{enumerate}[nosep]
  \item Lux acknowledges Liquidity.io as the canonical regulated DEX,
        ATS, broker--dealer (BD), and transfer agent (TA) counterparty
        for the digital-securities market.
  \item The Liquidity Protocol's OMA contract may be deployed as a
        first-class precompile-aware contract on Lux C-Chain (today)
        and on Lux D-Chain (when LP-134's D-Chain lands).
  \item Lux validators verify Liquidity-Protocol-compliant transaction
        envelopes as part of normal block validation; no special-case
        VM, no privileged signer.
  \item Liquidity-Protocol-compliant transactions inherit the same
        Quasar 3.0 finality predicate as native EVM transactions,
        including post-quantum finality, parallel-lane liveness, and
        subject-binding replay protection.
\end{enumerate}
This is a one-way fact about the protocol stack: Lux does not assume
operational dependency on the Liquidity.io platform; the protocol
remains an external service governed by Liquidity.io.
\end{definition}

\begin{remark}[Trade-name overlap]
The phrase ``liquidity protocol'' has prior usage in Zoo's NFT
liquidity literature (2021).  In Lux documents we always write
\emph{the Liquidity Protocol} (capitalised) for the 2026-04-01
Liquidity.io launch and \emph{NFT-LP} or \emph{Zoo NFT Liquidity
Protocol} for the 2021 Zoo construction.  Section~\ref{sec:no-collide}
formally separates the two.
\end{remark}

%% ==================================================================
\section{Threat and execution model}
\label{sec:model}
%% ==================================================================

The execution model is unchanged from Quasar 3.0:
\begin{itemize}[nosep]
  \item validator set $\mathcal{V}$ with stake $S$ and Byzantine
        fraction $f < 1/3$;
  \item three-lane finality certificate
        $\mathcal{C} = (\sigma_{\mathrm{BLS}}, \sigma_{\mathrm{RT}},
        \pi_{\mathrm{ZK}})$ over the certificate subject
        $\mathsf{subj}$ defined in~\cite{quasar-cert};
  \item parallel-execution layer QuasarSTM (LP-010, LP-135) with
        Block-STM serialisability and GPU residency~(LP-137).
\end{itemize}

The additional protocol surface is the Liquidity Protocol contract set
$\Pi_{\mathrm{LP}} = \{ \texttt{OMA}, \texttt{ATS}_\$,
\texttt{Oracle}, \texttt{ComplianceModule} \}$.  We treat
$\Pi_{\mathrm{LP}}$ as a closed sub-system whose behaviour is defined
by~\cite{lp-paper}.  A Liquidity-Protocol-compliant transaction is a
transaction whose call graph terminates in a successful invocation of
\texttt{OMA.swap} or \texttt{OMA.settle} after passing the protocol's
own pre-conditions (NBBO bound, inventory check, compliance gate).

\begin{definition}[LP-compliant transaction]
\label{def:lp-compliant}
A transaction $\mathsf{tx}$ is \emph{LP-compliant} on a Lux execution
chain $\mathcal{X} \in \{\text{C-Chain},\text{D-Chain}\}$ iff:
(a)~$\mathsf{tx}$ is a valid EVM transaction on $\mathcal{X}$;
(b)~$\mathsf{tx}$'s top-level call targets a contract in
$\Pi_{\mathrm{LP}}$;
(c)~the protocol-internal pre-conditions of~\cite{lp-paper} are met
(verified by the contract code itself);
(d)~the resulting state delta is recorded in the block whose subject
$\mathsf{subj}$ is finalised by a valid QuasarCert.
\end{definition}

%% ==================================================================
\section{Inheritance of Quasar 3.0 PQ finality}
\label{sec:inherit}
%% ==================================================================

The first adoption result is a direct corollary of QuasarCert
soundness.

\begin{theorem}[Inherited PQ finality]
\label{thm:inherited-finality}
Let $\mathsf{tx}$ be an LP-compliant transaction included in block $b$
of a Lux execution chain $\mathcal{X}$, and let $\mathcal{C}_b$ be a
valid QuasarCert over the subject $\mathsf{subj}_b$ of $b$.  Then for
every quantum PPT adversary $\mathcal{A}_Q$ controlling stake fraction
$<\tau$,
\[
  \Pr[\,\mathsf{tx} \text{ is reorged out of }
       \mathcal{X}\;|\;\mathcal{C}_b\text{ verifies}\,]
  \;\leq\;
  \mathsf{Adv}^{\mathrm{PQ\text{-}forge}}_{\mathcal{A}_Q}
  \;=\; \mathsf{negl}(\lambda).
\]
\end{theorem}

\begin{proof}
By Theorem~7.7 of~\cite{quasar-cert} (post-quantum safety of
QuasarCert), a forgery of the PQ predicate
$\mathsf{VerifyPQ} = \mathsf{RT.Verify} \land \mathsf{ZK.Verify}$
reduces to one of MLWE or MSIS.  Both are NIST-Level-3 PQ assumptions
under~\cite{nistpq}.  An LP-compliant $\mathsf{tx}$ is, by
Def.~\ref{def:lp-compliant}(a), a normal EVM transaction --- there is
no privileged path.  Reorging $\mathsf{tx}$ requires forging an
alternative $\mathcal{C}_b'$ on a competing subject
$\mathsf{subj}_b' \neq \mathsf{subj}_b$, which by Theorem~7.7
of~\cite{quasar-cert} is bounded by
$\mathsf{Adv}^{\mathrm{PQ\text{-}forge}}_{\mathcal{A}_Q} = \mathsf{negl}(\lambda)$.
The Liquidity Protocol contract layer adds new \emph{state},
not new \emph{certificates}; the cert lane is unchanged.
\end{proof}

\begin{corollary}[Settlement irrevocability]
\label{cor:settlement-irrevocable}
Once an OMA settlement is included in a Lux block whose QuasarCert
verifies, the settlement is irrevocable except with negligible
probability under PQ assumptions, even against quantum adversaries
running Shor and Grover on classical signing components.
\end{corollary}

%% ==================================================================
\section{Atomicity composes with QuasarSTM serialisability}
\label{sec:compose}
%% ==================================================================

QuasarSTM (LP-010, LP-135) provides Block-STM-style optimistic
concurrency control with serialisability as the safety
property~\cite{lp135}.  The Liquidity Protocol's settlement contract
guarantees \emph{trade atomicity}: a swap either fully executes
(token in / token out / inventory updated) or fully reverts.  We
prove the two compose.

\begin{definition}[Trade atomicity, LP \S5.2]
\label{def:lp-atomic}
A Liquidity-Protocol swap from token $T_a$ for token $T_b$ at
oracle price $P_o$ on chain $\mathcal{X}$ either executes the joint
state update
\[
  s' = s
       \oplus (\Delta T_a := -\amt,\;
               \Delta T_b := +\amt \cdot P_o \cdot (1-\spread),\;
               \Delta\bal_{\mathrm{ATS}} := -\amt \cdot P_o)
\]
atomically inside a single EVM transaction, or it reverts and
$s' = s$.
\end{definition}

\begin{lemma}[Block-STM commutes with single-tx atomicity]
\label{lem:bstm-atomic}
Let $\mathsf{tx}_{\mathrm{LP}}$ be an LP-compliant transaction with
trade atomicity (Def.~\ref{def:lp-atomic}).  Block-STM execution of
$\mathsf{tx}_{\mathrm{LP}}$ inside a Lux block produces the same
output state as a serial execution of the same block.
\end{lemma}

\begin{proof}
Block-STM~\cite{blockstm} executes transactions optimistically and
re-executes any whose read--write set conflicts with a higher-priority
sibling.  A re-execution simply re-runs the EVM call graph from
scratch; the Liquidity Protocol contract's state machine is
deterministic given block-level inputs (oracle round, ATS balance,
compliance gate).  Therefore re-executions converge.  Atomicity is a
\emph{single-transaction} property; Block-STM never partially commits
a transaction (\cite[Lemma 4.3]{blockstm}).  Hence the trade-atomicity
invariant is preserved under Block-STM scheduling, and the final
serialised order is some valid total order over the block's transactions.
\end{proof}

\begin{theorem}[Composition of LP atomicity and QuasarSTM serialisability]
\label{thm:compose}
For any Lux block $b$ on a QuasarSTM-enabled execution chain
$\mathcal{X}$ containing LP-compliant transactions
$T_{\mathrm{LP}} = \{\mathsf{tx}_1, \ldots, \mathsf{tx}_m\}$
interleaved with arbitrary EVM transactions $T_{\mathrm{EVM}}$, the
final state $s_b'$ satisfies:
\begin{enumerate}[nosep]
  \item every $\mathsf{tx}_i \in T_{\mathrm{LP}}$ is either fully
        committed (atomic) or fully reverted in $s_b'$;
  \item there exists a serial schedule $\pi$ of
        $T_{\mathrm{LP}} \cup T_{\mathrm{EVM}}$ such that running
        $\pi$ over the pre-state $s_{b-1}$ produces $s_b'$;
  \item $s_b'$ is unique, deterministic, and finalised by the
        QuasarCert $\mathcal{C}_b$.
\end{enumerate}
\end{theorem}

\begin{proof}
(1) is Lemma~\ref{lem:bstm-atomic}.  (2) is Block-STM serialisability
(LP-135 \S6, machine-checked in
\texttt{lean/Consensus/DAGEVM.lean}~\cite{dagevm-lean}).
(3)~combines (2) with QuasarCert subject-binding: the unique post-state
$s_b'$ commits to the C-Chain (or D-Chain) state root inside
$\mathsf{subj}_b$, which the QuasarCert finalises by
Theorem~7.7 of~\cite{quasar-cert}.
\end{proof}

\begin{corollary}[No LP-induced contention with QuasarSTM]
\label{cor:no-contention}
Adoption of the Liquidity Protocol does not weaken QuasarSTM 4.0
(LP-135) properties: serialisability still holds; throughput bounds
in LP-135 \S8 are unchanged; the only contention is intra-OMA
between concurrent swaps on the same token pair, which Block-STM
handles by re-execution.
\end{corollary}

%% ==================================================================
\section{Cross-chain replay protection}
\label{sec:replay}
%% ==================================================================

The Liquidity Protocol may deploy on multiple EVM chains
simultaneously (Liquid EVM, Lux C-Chain, partner L1s).  An adversary
could attempt to replay a Lux-finalised OMA settlement on a different
chain.  We show this is precluded by QuasarCert subject-binding.

\begin{theorem}[Cross-chain replay is hard]
\label{thm:replay}
Let $\mathsf{tx}$ be an LP-compliant transaction finalised on Lux
execution chain $\mathcal{X}_1$ with QuasarCert $\mathcal{C}_1$ over
subject $\mathsf{subj}_1$, and let $\mathcal{X}_2 \neq \mathcal{X}_1$
be any other Lux execution chain.  An adversary $\mathcal{A}_Q$ that
constructs a valid QuasarCert $\mathcal{C}_2$ over a subject
$\mathsf{subj}_2$ that re-executes $\mathsf{tx}$ on $\mathcal{X}_2$
must have advantage at most
$\mathsf{Adv}^{\mathrm{PQ\text{-}forge}}_{\mathcal{A}_Q} =
\mathsf{negl}(\lambda)$.
\end{theorem}

\begin{proof}
By Def.~6.1 of LP-105~\cite{quasar-cert} the certificate subject
$\mathsf{subj}$ commits to \texttt{chain\_id}.  Therefore
$\mathsf{subj}_1 \neq \mathsf{subj}_2$ as soon as
$\mathcal{X}_1 \neq \mathcal{X}_2$.  Any adversary forging
$\mathcal{C}_2$ violates Theorem~7.8 of~\cite{quasar-cert}
(subject-binding replay protection).  By Theorem~7.7 the bound is
$\mathsf{negl}(\lambda)$ even against quantum adversaries.
\end{proof}

\begin{remark}[Liquid EVM]
The Liquidity Network's Liquid EVM (chain id 8675309 / 8675310 /
8675311) is operationally distinct from Lux C-Chain.  This theorem
addresses replay between two \emph{Lux} execution chains
(C-Chain and the forthcoming D-Chain).  Replay between Liquid EVM
and Lux is a separate cross-network concern handled by
\cite{lp-paper} \S7.
\end{remark}

%% ==================================================================
\section{No collision with Zoo NFT-LP}
\label{sec:no-collide}
%% ==================================================================

The 2021 Zoo whitepaper~\cite{zoo-2021} coined the term
\emph{NFT Liquidity Protocol} for collateral-backed NFT lending and
sustainability-funded yield.  The 2026 Liquidity.io launch uses the
same English noun phrase for a securities-collateralised
oracle-mirrored AMM.  The two are mathematically distinct.

\begin{proposition}[Disjoint asset classes]
\label{prop:disjoint}
The state spaces $\mathcal{S}_{\mathrm{Zoo\text{-}NFT\text{-}LP}}$
and $\mathcal{S}_{\mathrm{Liquidity\text{-}LP}}$ are disjoint:
\[
  \mathcal{S}_{\mathrm{Zoo\text{-}NFT\text{-}LP}}
  \subseteq \{\,\text{ERC-721}\,\}^* \times
            \{\,\text{ERC-20 collateral}\,\}^*,
  \qquad
  \mathcal{S}_{\mathrm{Liquidity\text{-}LP}}
  \subseteq \{\,\text{ERC-3643 securities}\,\}^* \times
            \{\,\text{USDL stable}\,\}^*.
\]
A Zoo NFT-LP transaction can never be mistaken for a Liquidity
Protocol transaction (different ABI selectors, different contract
addresses, different token standards) at the EVM dispatch layer.
\end{proposition}

\begin{proof}
By inspection of contract bytecode: Zoo NFT-LP collateralises
\texttt{ERC-721}; Liquidity Protocol settles \texttt{ERC-3643}.  The
EVM call selector is a fixed 4-byte function discriminator; the
selectors differ.  The disjointness is structural, not nominal.
\end{proof}

Composition between the two systems is therefore safe: a single Lux
execution chain can host both.  Zoo's ZIP-0206 NFT-LP proofs are not
affected by Liquidity Protocol adoption, and vice versa.  See
\cite{zoo-nft-lp} for the formal Zoo proof.

%% ==================================================================
\section{Adoption record}
\label{sec:record}
%% ==================================================================

\begin{table}[h]
\centering
\caption{Lux Network protocol-stack chronology, 2025--2026}
\begin{tabular}{lll}
\toprule
Date & Event & Reference \\
\midrule
2025-12-25 & Quasar 3.0 chain activation
           & LP-020, \cite{quasar-cert} \\
2026-02-14 & QuasarSTM 4.0 production activation
           & LP-135 \\
2026-04-01 & Liquidity Protocol launch (Liquidity.io)
           & \cite{lp-paper} \\
2026-04-20 & \textbf{Lux adopts the Liquidity Protocol}
           & \emph{this paper} \\
\bottomrule
\end{tabular}
\end{table}

\paragraph{Future work (post 2026-04-20).}
LP-134 reserves D-Chain for the Liquidity stack; D-Chain bring-up,
GPU-residency invariants for the OMA contract under LP-137, and
extended cross-chain composition with M-Chain MPC are deferred.

%% ==================================================================
\section{Worked example: a regulated swap}
\label{sec:example}
%% ==================================================================

To make Theorem~\ref{thm:compose} concrete, consider a user $U$
holding 100 USDL on Lux C-Chain who wishes to acquire AAPL
(NYSE security mirrored as ERC-3643 on Liquidity infrastructure)
via the Liquidity Protocol's OMA. The transaction proceeds:

\begin{enumerate}[nosep]
  \item $U$ submits $\mathsf{tx} = \texttt{OMA.swap(USDL, AAPL, 100)}$
        to a Lux validator.
  \item The validator executes $\mathsf{tx}$ inside the C-Chain EVM
        under QuasarSTM 4.0:
        \begin{itemize}[nosep]
          \item read oracle: $P_o(\text{AAPL}) = \$237.41$
                (LP \S5.1 oracle interface);
          \item read ATS account inventory:
                $\bal_{\mathrm{ATS}}(\text{AAPL}) = 50{,}000$ shares;
          \item compliance gate: ERC-3643 modules check
                $U$'s KYC / accreditation status (BD service);
          \item if pass, transfer $100$ USDL out, mint
                $0.421$ AAPL ERC-3643 tokens to $U$, debit ATS
                inventory.
        \end{itemize}
  \item Block $b$ commits with all such transactions; QuasarCert
        $\mathcal{C}_b$ over $\mathsf{subj}_b$ is signed by the
        validator quorum.
  \item Quasar 3.0 finalises $b$.
\end{enumerate}

By Theorem~\ref{thm:inherited-finality}, after $\mathcal{C}_b$
verifies, $U$'s 0.421 AAPL position is finalised under PQ
guarantees: even a quantum adversary running Shor on the BLS lane
cannot revert the swap, because revert requires forging the PQ lanes
$\sigma_{\mathrm{RT}} \land \pi_{\mathrm{ZK}}$
(Theorem~7.7 of \cite{quasar-cert}).

By Theorem~\ref{thm:compose}, the swap commits atomically: $U$
either has 0.421 AAPL (and $-100$ USDL), or has full original
balance --- never partial.

By Theorem~\ref{thm:replay}, the swap cannot be replayed on
D-Chain (or any other Lux execution chain) without forging a fresh
QuasarCert with the alternative chain id, which is equally hard.

%% ==================================================================
\section{Concrete parameter analysis}
\label{sec:params}
%% ==================================================================

\begin{table}[h]
\centering
\caption{Cryptographic parameters and security levels (NIST PQ
classification).}
\begin{tabular}{llll}
\toprule
Component & Standard & Hardness & NIST Level \\
\midrule
BLS12-381 aggregate & RFC 9380     & co-CDH (classical)  & --- (no PQ)\\
Ringtail              & LP-105       & MLWE + MSIS         & 3 \\
ML-DSA-65             & FIPS 204     & MLWE + MSIS         & 3 \\
ML-KEM-768            & FIPS 203     & MLWE                & 3 \\
SLH-DSA-256           & FIPS 205     & second-preimage     & 5 \\
Groth16 (Z-Chain)     & --           & KEA / AGM           & --- \\
\bottomrule
\end{tabular}
\end{table}

The PQ-only bound (Theorem 7.7 of \cite{quasar-cert}) requires
$\mathsf{VerifyPQ} = \mathsf{RT.Verify} \land \mathsf{ZK.Verify}$.
Both lanes operate at NIST Level 3 (Ringtail Module-LWE) and are
auxiliary-bound (Groth16 KEA + per-signature ML-DSA).
Concrete bit-security $\geq 192$ bits classical, $\geq 100$ bits
quantum (Grover-corrected).

%% ==================================================================
\section{Implementation footprint on Lux}
\label{sec:impl}
%% ==================================================================

\begin{itemize}[nosep]
  \item No new VM: the Liquidity Protocol contracts compile to
        standard EVM bytecode and deploy on C-Chain.
  \item No special-case validator code: validators verify
        Liquidity-Protocol transactions identically to any other
        EVM transaction.
  \item Oracle / OMA / Compliance contracts are deployed at
        published addresses; their ABI is part of the
        Liquidity.io spec (\cite{lp-paper} \S4).
  \item Gas: per-swap gas budget $\leq 250{,}000$ gas
        (\cite{lp-paper} \S8); fits comfortably in a 30M-gas Lux
        block.
\end{itemize}

The adoption is therefore a contract-layer fact about deployment
practice, not a protocol fork or hard-fork. No Lux node software
needs to update for adoption to take effect.

%% ==================================================================
\begin{thebibliography}{9}

\bibitem{quasar-cert}
Z.~Kelling.
\newblock QuasarCert Soundness: Canonical Formal Specification, Quasar 3.0.
\newblock \texttt{luxfi/proofs/quasar-cert-soundness.tex}, 2025-12-15.

\bibitem{lp-paper}
Liquidity.io.
\newblock The Liquidity Protocol: Oracle-Mirrored AMM for Universal
         Cross-Chain Liquidity, v1.0.
\newblock \texttt{liquidityio/proofs/liquidity-protocol/}, April~2026.

\bibitem{lp134}
Lux Industries.
\newblock LP-134 Lux Chain Topology.
\newblock \texttt{luxfi/lps/LP-134-lux-chain-topology.md}, 2025-12-15.

\bibitem{lp135}
Lux Industries.
\newblock LP-135 QuasarSTM 4.0 Production Spec.
\newblock \texttt{luxfi/lps/LP-135-quasarstm-4-research.md}, 2025-12-25.

\bibitem{nistpq}
NIST.
\newblock FIPS 203 / 204 / 205: Module-Lattice and Hash-Based Standards.
\newblock 2024.

\bibitem{blockstm}
R.~Gelashvili et al.
\newblock Block-STM: Scaling Blockchain Execution by Turning Ordering
         Curse to a Performance Blessing.
\newblock PPoPP 2023.

\bibitem{dagevm-lean}
Lux Industries.
\newblock \texttt{lean/Consensus/DAGEVM.lean}: machine-checked
         Block-STM-as-consensus serialisability.
\newblock \texttt{luxfi/proofs/lean/}.

\bibitem{zoo-2021}
A.~Worring, Z.~Kelling.
\newblock Zoo: NFT-Driven Conservation, October 2021.
\newblock \texttt{zooai/papers/zoo-2021-original-whitepaper/}.

\bibitem{zoo-nft-lp}
Zoo Labs Foundation.
\newblock Zoo NFT Liquidity Protocol: Soundness.
\newblock \texttt{zooai/proofs/zoo-nft-liquidity-protocol-soundness.tex}, 2026.

\end{thebibliography}

\end{document}
